%% Supplementary material for:
%% "Day-ahead probabilistic forecasting of residential rooftop photovoltaic power:
%%  transfer to unseen systems and the input-imposed accuracy limit"
%% Compile separately from main.tex. Uses the same refs.bib.

\documentclass[11pt]{article}
\usepackage[a4paper,margin=25mm]{geometry}
\usepackage{amsmath}
\usepackage{booktabs}
\usepackage{tabularx}
\usepackage[authoryear,round]{natbib}
\usepackage[hidelinks]{hyperref}
\usepackage{url} % \path{} lets long file paths break at / _ .
\def\UrlBreaks{\do\/\do\_\do\.}

\renewcommand{\thetable}{S\arabic{table}}
\renewcommand{\thesection}{S\arabic{section}}

\title{Supplementary material\\[4pt]\large Day-ahead probabilistic forecasting of residential rooftop photovoltaic power: transfer to unseen systems and the input-imposed accuracy limit}
\author{Chandradip Roy, Sara Ghane}
\date{}

\begin{document}
\maketitle

Section and table numbers in the main text are referred to as ``main text''. All values come from the analysis repository (tag \texttt{v1.1-paper}); the source file is given under each table.

%% ---------------------------------------------------------------------------
\section{Headline arms and recalibration}
\label{sec:s-arms}

\begin{table}[!htb]
\centering
\caption{Headline arms on the chronological test split, before and after conformal recalibration. Systems: 33 = full fleet, 32 = system 81376 excluded. Source: \path{models/consolidated/exp4_arms.csv}.}
\label{tab:s-arms}
\small
\begin{tabularx}{\textwidth}{@{}l X c c c c c@{}}
\toprule
Arm & Configuration & Systems & $R^2$ & Skill & Pinball & P25--P90 cov. \\
\midrule
A & Base (12 features) & 33 & 0.7960 & 0.4139 & 0.02606 & 0.619 \\
  & \quad recalibrated & 33 & 0.7964 & 0.4145 & 0.02602 & 0.658 \\
B & Base + multi-model mean (proposed) & 33 & 0.7998 & 0.4194 & 0.02580 & 0.610 \\
  & \quad recalibrated & 33 & 0.8000 & 0.4197 & 0.02578 & 0.646 \\
D & Base (12 features) & 32 & 0.8070 & 0.4285 & 0.02558 & 0.622 \\
  & \quad recalibrated & 32 & 0.8074 & 0.4292 & 0.02554 & 0.659 \\
C & Base + multi-model mean & 32 & 0.8114 & 0.4351 & 0.02529 & 0.614 \\
  & \quad recalibrated & 32 & 0.8117 & 0.4355 & 0.02527 & 0.649 \\
\bottomrule
\end{tabularx}
\end{table}

Arms C and D exclude system 81376 as a sensitivity analysis. Its defect is time-varying and was found after the test errors were seen. Its April--May output relative to the fleet median was 0.85 in 2024 and 0.89 in 2025, but 0.60 in 2026 (test window), the largest drop in the fleet (fleet median change $-3\%$). Source: \path{src/eval/system_drift.py}, \path{models/system_drift_monthly.csv}.

%% ---------------------------------------------------------------------------
\section{Selected hyperparameters}
\label{sec:s-hparams}

\begin{table}[!htb]
\centering
\caption{LightGBM hyperparameters. Temporal: one search per quantile on the chronological split. Transfer: P50 search against held-out systems. Sources: \path{models/tuned/best_params_quantile.json}, \texttt{best\_params\_transfer.json}.}
\label{tab:s-hparams}
\small
\begin{tabular}{@{}l c c c c@{}}
\toprule
Parameter & Temporal P25 & Temporal P50 & Temporal P90 & Transfer P50 \\
\midrule
Learning rate & 0.0152 & 0.0129 & 0.0147 & 0.0184 \\
Trees & 1418 & 1999 & 1993 & 1500 \\
Max depth & 11 & 15 & 20 & 4 \\
Leaves & 222 & 247 & 148 & 101 \\
Min child samples & 45 & 155 & 82 & 119 \\
Subsample (freq.) & 0.771 (1) & 0.630 (0) & 0.779 (4) & 0.561 (2) \\
Column subsample & 0.631 & 0.590 & 0.567 & 0.668 \\
L1 penalty & 0.122 & 0.163 & 0.492 & 0.144 \\
L2 penalty & 0.0347 & 0.0081 & 0.0015 & 4.344 \\
\bottomrule
\end{tabular}
\end{table}

Search spaces: learning rate 0.01--0.2 (log scale); tree caps 2000 (temporal) and 1500 (transfer). Both searches reached or approached their caps (main text, Limitations).

%% ---------------------------------------------------------------------------
\section{Irradiance-level MOS}
\label{sec:s-mos}

\begin{table}[!htb]
\centering
\caption{MOS of ECMWF day-ahead GHI on the test split ($n = 37{,}012$ hourly rows), by truth source and learner. MBE = forecast minus truth. Positive reduction = improvement. Source: \path{models/mos_satellite_metrics.csv}.}
\label{tab:s-mos}
\small
\begin{tabular}{@{}l l c c c c c@{}}
\toprule
Truth & Model & MAE & RMSE & MBE & $R^2$ & MAE red. (\%) \\
 & & (W\,m$^{-2}$) & (W\,m$^{-2}$) & (W\,m$^{-2}$) & & \\
\midrule
SARAH-3 & Raw ECMWF & 55.49 & 86.27 & +11.40 & 0.866 & 0.00 \\
 & Ridge & 55.19 & 83.14 & +0.88 & 0.876 & 0.54 \\
 & Random forest & 51.96 & 81.94 & +3.73 & 0.879 & 6.37 \\
 & LightGBM & 55.82 & 86.36 & +2.15 & 0.866 & $-0.60$ \\
 & XGBoost & 56.55 & 86.48 & +0.93 & 0.866 & $-1.91$ \\
\midrule
ERA5 & Raw ECMWF & 24.51 & 40.94 & +0.81 & 0.968 & 0.00 \\
 & Ridge & 24.74 & 41.14 & +1.13 & 0.968 & $-0.96$ \\
 & Random forest & 24.81 & 42.14 & +0.42 & 0.966 & $-1.20$ \\
 & LightGBM & 25.60 & 43.63 & +0.92 & 0.964 & $-4.46$ \\
 & XGBoost & 26.85 & 45.57 & $-0.72$ & 0.961 & $-9.56$ \\
\bottomrule
\end{tabular}
\end{table}

The test rows here are a subset of the study window, so the raw errors differ from the whole-window reference comparison in the main text (RMSE 86.6 against 51.1~W\,m$^{-2}$). Both comparisons point the same way. Whole-window source: \path{src/eval/reference_comparison.py}.

%% ---------------------------------------------------------------------------
\section{Local calibration sweep}
\label{sec:s-calib}

\begin{table}[!htb]
\centering
\caption{Median skill on 33 held-out rooftops by calibration window and correction. The evaluation window shrinks as the calibration window grows, so compare only within a row. Source: \path{models/consolidated/exp5_calib_skill_pivot.csv}.}
\label{tab:s-calib-skill}
\small
\begin{tabular}{@{}r c c c c@{}}
\toprule
Days & None & Bias & Capacity + bias & Linear \\
\midrule
0 & 0.3465 & -- & -- & -- \\
3 & \textbf{0.3469} & 0.2917 & 0.2625 & $-0.1099$ \\
7 & \textbf{0.3456} & 0.3260 & 0.3062 & $-0.1122$ \\
14 & \textbf{0.3404} & 0.3069 & 0.3151 & 0.1169 \\
30 & 0.3186 & 0.3198 & 0.3230 & \textbf{0.3276} \\
45 & 0.3095 & 0.3299 & 0.3192 & \textbf{0.3312} \\
60 & 0.2660 & 0.2912 & 0.2809 & \textbf{0.3194} \\
\bottomrule
\end{tabular}
\end{table}

\begin{table}[!htb]
\centering
\caption{As Table~\ref{tab:s-calib-skill}, median $R^2$. Source: \path{models/consolidated/exp5_calib_r2_pivot.csv}.}
\label{tab:s-calib-r2}
\small
\begin{tabular}{@{}r c c c c@{}}
\toprule
Days & None & Bias & Capacity + bias & Linear \\
\midrule
0 & 0.7043 & -- & -- & -- \\
3 & 0.7100 & 0.7019 & 0.6871 & 0.2309 \\
7 & 0.7185 & 0.7209 & 0.7074 & 0.3322 \\
14 & 0.7259 & 0.7116 & 0.7203 & 0.5424 \\
30 & 0.7366 & 0.7390 & 0.7441 & 0.7366 \\
45 & 0.7418 & 0.7532 & 0.7545 & 0.7553 \\
60 & 0.7408 & 0.7626 & 0.7584 & 0.7551 \\
\bottomrule
\end{tabular}
\end{table}

By $R^2$, the best correction from 30 days onward varies between methods; by skill, the main-text criterion, linear rescaling is best at every window of 30 days or more.

%% ---------------------------------------------------------------------------
\section{Shading flag and zero-shot accuracy}
\label{sec:s-shade}

Owners report shading as none, low or medium. Median zero-shot $R^2$ falls across these categories: 0.748 (none, 26 systems), 0.658 (low, 3) and 0.643 (medium, 3); one system has no entry. With three systems in each shaded category the trend is suggestive only. Sources: \path{models/consolidated/exp5_loso_per_system.csv}, \path{data/interim/pvoutput_selection/systems_metadata.csv}.

%% ---------------------------------------------------------------------------
\section{Paired comparisons}
\label{sec:s-paired}

\begin{table}[!htb]
\centering
\caption{Paired differences (second arm minus first) on the chronological test split. Intervals: 95\% day-block bootstrap, 2000 resamples of 107 test days. $p$: two-sided Diebold--Mariano test on daily squared error of P50 ($p_\mathrm{SE}$) and daily mean pinball loss ($p_\mathrm{pin}$), Newey--West variance. Pinball differences in units of $10^{-3}$. Source: \path{src/eval/revision_stats.py}, \path{models/revision/paired_comparisons.csv}.}
\label{tab:s-paired}
\scriptsize
\begin{tabularx}{\textwidth}{@{}X c c c c@{}}
\toprule
Comparison & $\Delta$ skill & $\Delta$ pinball & $p_\mathrm{SE}$ & $p_\mathrm{pin}$ \\
\midrule
E0 objective: pinball vs RMSE selection (33 systems) & $+0.0031$\,[$+0.0018$, $+0.0045$] & $-0.20$\,[$-0.24$, $-0.16$] & $<0.001$ & $<0.001$ \\
E1 oracle: + regional truth vs NWP only & $+0.0341$\,[$+0.0224$, $+0.0462$] & $-1.51$\,[$-1.99$, $-1.04$] & $<0.001$ & $<0.001$ \\
E1 oracle: regional truth without NWP vs NWP only & $-0.1211$\,[$-0.1544$, $-0.0891$] & $+7.53$\,[$+5.89$, $+9.06$] & $<0.001$ & $<0.001$ \\
E3: + multi-model mean vs base (33 systems) & $+0.0055$\,[$+0.0001$, $+0.0106$] & $-0.26$\,[$-0.46$, $-0.06$] & 0.04 & 0.01 \\
E3: + mean of all 12 members vs base & $+0.0004$\,[$-0.0035$, $+0.0043$] & $-0.03$\,[$-0.19$, $+0.12$] & 0.96 & 0.77 \\
E3: + ensemble-mean POA vs base & $+0.0008$\,[$-0.0039$, $+0.0054$] & $-0.01$\,[$-0.20$, $+0.18$] & 0.81 & 0.96 \\
E3: + lagged mean (days 1-3) vs base & $+0.0002$\,[$-0.0020$, $+0.0022$] & $+0.03$\,[$-0.08$, $+0.14$] & 0.88 & 0.35 \\
E3: + spread only vs base & $-0.0009$\,[$-0.0042$, $+0.0020$] & $+0.43$\,[$+0.24$, $+0.63$] & 0.47 & $<0.001$ \\
E3: + all 10 ensemble features vs base & $+0.0016$\,[$-0.0046$, $+0.0071$] & $+0.30$\,[$+0.04$, $+0.60$] & 0.90 & 0.03 \\
E4: recalibrated vs raw (33 systems, base + ens) & $+0.0003$\,[$-0.0003$, $+0.0009$] & $-0.02$\,[$-0.08$, $+0.04$] & 0.30 & 0.64 \\
E4 (32 systems): + multi-model mean vs base & $+0.0065$\,[$+0.0011$, $+0.0118$] & $-0.29$\,[$-0.50$, $-0.09$] & 0.02 & 0.005 \\
MOS seen rooftops (33): corrected irradiance vs baseline & $-0.0004$\,[$-0.0021$, $+0.0013$] & $+0.12$\,[$+0.03$, $+0.23$] & 0.46 & 0.04 \\
MOS seen rooftops (33): MOS as extra features vs baseline & $+0.0000$\,[$-0.0022$, $+0.0023$] & $+0.18$\,[$+0.07$, $+0.32$] & 0.94 & 0.01 \\
MOS seen rooftops (32): corrected irradiance vs baseline & $-0.0005$\,[$-0.0021$, $+0.0011$] & $+0.10$\,[$+0.02$, $+0.20$] & 0.42 & 0.06 \\
\bottomrule
\end{tabularx}
\end{table}

%% ---------------------------------------------------------------------------
\section{Seed variation of the zero-shot results}
\label{sec:s-seeds}

\begin{table}[!htb]
\centering
\caption{Leave-one-system-out P50 results on 33 held-out rooftops for each random seed (LightGBM row and column subsampling). Failures: skill $\le 0$. Paired Wilcoxon signed-rank tests on seed-averaged per-rooftop $R^2$: transfer better than temporal on 22 of 33 rooftops ($p = 0.001$), transfer better than preliminary on 19 ($p = 0.015$), temporal better than preliminary on 10 ($p = 0.014$). Source: \path{src/models/loso_seeds.py}, \path{models/revision/loso_seeds_summary.csv}, \path{src/eval/revision_stats.py}.}
\label{tab:s-seeds}
\small
\begin{tabularx}{\textwidth}{@{}X c c c c c c c@{}}
\toprule
Configuration & Seed & Median $R^2$ & Mean $R^2$ & P10 & Worst & Median skill & Failures \\
\midrule
Preliminary & 0 & 0.7195 & 0.5967 & 0.3622 & $-0.8550$ & 0.3259 & 4 \\
Preliminary & 1 & 0.7239 & 0.6040 & 0.3808 & $-0.7164$ & 0.3046 & 4 \\
Preliminary & 2 & 0.7153 & 0.5997 & 0.3603 & $-0.7696$ & 0.2987 & 4 \\
Temporal & 0 & 0.7058 & 0.5856 & 0.3113 & $-0.9989$ & 0.3111 & 5 \\
Temporal & 1 & 0.7051 & 0.5875 & 0.2958 & $-0.8876$ & 0.3075 & 5 \\
Temporal & 2 & 0.7103 & 0.5911 & 0.3288 & $-0.9205$ & 0.3179 & 4 \\
Transfer & 0 & 0.7093 & 0.6576 & 0.4978 & $-0.1565$ & 0.3457 & 2 \\
Transfer & 1 & 0.7046 & 0.6543 & 0.4961 & $-0.1431$ & 0.3403 & 2 \\
Transfer & 2 & 0.7072 & 0.6584 & 0.5023 & $-0.1526$ & 0.3411 & 2 \\
Transfer & 3 & 0.7172 & 0.6577 & 0.4974 & $-0.1489$ & 0.3446 & 2 \\
Transfer & 4 & 0.7163 & 0.6594 & 0.5003 & $-0.1580$ & 0.3418 & 2 \\
\bottomrule
\end{tabularx}
\end{table}

%% ---------------------------------------------------------------------------
\section{Reference baselines}
\label{sec:s-baselines}

\begin{table}[!htb]
\centering
\caption{Baselines without machine learning on the chronological test split, 33 systems. Pooled: over all test rows. Per rooftop: distribution of $R^2$ over the 33 systems; failures, skill $\le 0$. Physical chain: Hay--Davies plane-of-array irradiance from the ECMWF forecast, Faiman cell temperature, PVWatts ($\gamma = -0.4\,\%\,\mathrm{K}^{-1}$), 14\% losses, nameplate capacity. Source: \path{src/models/baselines.py}, \path{models/revision/baselines_pooled.csv}, \texttt{baselines\_per\_system.csv}.}
\label{tab:s-baselines}
\small
\begin{tabularx}{\textwidth}{@{}X c c c c c c@{}}
\toprule
 & \multicolumn{2}{c}{Pooled} & \multicolumn{4}{c}{Per rooftop} \\
\cmidrule(lr){2-3}\cmidrule(l){4-7}
Baseline & $R^2$ & Skill & Median $R^2$ & P10 & Worst & Failures \\
\midrule
Persistence & 0.4062 & 0 & 0.409 & 0.267 & 0.055 & -- \\
Smart persistence & 0.4006 & $-0.0047$ & 0.404 & 0.260 & 0.045 & 33 \\
Physical chain, nameplate & 0.5205 & 0.1014 & 0.608 & 0.160 & $-1.445$ & 8 \\
Physical chain, per-rooftop scale & 0.6316 & 0.2124 & 0.633 & 0.405 & 0.134 & 2 \\
\bottomrule
\end{tabularx}
\end{table}

Smart persistence rescales yesterday's output by the ratio of clear-sky GHI between the two days. At a one-day lag that ratio is close to one, so the two persistence references differ by less than 0.5\% of RMSE. The per-rooftop scale factor is fitted by least squares on each rooftop's training split (median 0.882, range 0.494--1.584). Across the 33 rooftops, the bias of the nameplate physical chain correlates with the capacity ratio $P_{99}/C_\mathrm{p}$ at +0.845, and its $R^2$ correlates with the zero-shot learner's $R^2$ at +0.938. Wind speed from Open-Meteo is served in km\,h$^{-1}$ and is converted to m\,s$^{-1}$; plane-of-array irradiance is capped at 1400~W\,m$^{-2}$ to remove rare transposition spikes at low sun.

%% ---------------------------------------------------------------------------
%% [EPSR-SCOPE] grid-facing aggregation. Optional for other journals (JOURNAL_NOTES.md).
\section{Aggregated forecast error}
\label{sec:s-aggregate}

\begin{table}[!htb]
\centering
\caption{Day-ahead P50 error of the proposed configuration (33 systems) after summing rooftop forecasts in kW. nRMSE: RMSE over installed capacity of the summed rooftops. Skill: against the summed persistence forecast. Daily imbalance: sum over days of $|\sum_t (P - \hat{P})|\,\Delta t$ as a share of energy. Individual: median over rooftops. Source: \path{src/eval/grid_aggregate.py}, \path{models/revision/grid_aggregate.csv}.}
\label{tab:s-aggregate}
\small
\begin{tabularx}{\textwidth}{@{}X c c c c c@{}}
\toprule
Level & Rooftops & nRMSE (\% cap.) & Skill & Daily imbalance (\%) & Mean daily imbalance (kWh) \\
\midrule
Individual rooftop (median) & 1 & 11.09 & 0.442 & 11.86 & 2.07 \\
Region: Ardennes & 4 & 8.62 & 0.476 & 10.22 & 8.10 \\
Region: Li\`ege & 6 & 8.52 & 0.518 & 10.13 & 11.06 \\
Region: West Flanders coast & 7 & 7.76 & 0.492 & 12.22 & 13.79 \\
Region: Antwerp & 8 & 7.70 & 0.481 & 11.46 & 14.97 \\
Region: Brussels and centre & 8 & 6.64 & 0.529 & 8.86 & 14.33 \\
Fleet & 33 & 4.99 & 0.590 & 8.42 & 49.99 \\
\bottomrule
\end{tabularx}
\end{table}

%% ---------------------------------------------------------------------------
\section{Positioning against published work}
\label{sec:s-position}

\begin{table}[!htb]
\centering
\caption{Comparison with published results. Validation: should agree if the pipeline is correct. Corroboration: qualitative agreement under different conditions. Not comparable: different quantity or reference.}
\label{tab:s-position}
\small
\begin{tabularx}{\textwidth}{@{}>{\raggedright\arraybackslash}p{2.6cm} >{\raggedright\arraybackslash}X >{\raggedright\arraybackslash}X@{}}
\toprule
Kind & Published & This study \\
\midrule
Validation: ECMWF GHI bias & +8~W\,m$^{-2}$, four years of ceilometer and pyranometer data at 60$^\circ$N \citep{Tuononen2019} & +9.2~W\,m$^{-2}$ against SARAH-3 \\
Validation: reference uncertainty & CM SAF satellite products, MAD 8--13~W\,m$^{-2}$ in central Europe \citep{Urraca2017} & inherited by all irradiance figures \\
Corroboration: ERA5 as reference & cloud-dependent bias, larger errors than satellite products \citep{Urraca2018} & ERA5 understates ECMWF RMSE 1.69-fold and reverses the bias \\
Corroboration: learner & tree ensembles best on 152 distributed systems \citep{Visser2019} & gradient boosting adopted \\
Corroboration: input vs model & NWP worth +11.6 points of skill at day-ahead, above any model family \citep{Nguyen2026} & oracle headroom +0.023 $R^2$; best input widening +0.0038 $R^2$ \\
Corroboration: transfer & local model needs $\sim$2 months to beat a naive baseline \citep{Depoortere2025} & local correction harmful below $\sim$30 days \\
Not comparable: probabilistic & 18.4\% mean-normalised CRPS, 14 utility-scale plants \citep{MayerYang2023} & mean pinball 0.02580; nRMSE 11.3\% of capacity \\
Not comparable: skill & skill against smart persistence \citep{Visser2019} & skill 0.4194 against plain persistence (smart persistence: $-0.005$) \\
\bottomrule
\end{tabularx}
\end{table}

%% ---------------------------------------------------------------------------
\section{Computation}
\label{sec:s-compute}

All experiments ran on one workstation with 12 logical CPU cores (Python 3.12, LightGBM 4.6.0, scikit-learn 1.9.0, pandas 3.0.3, pvlib). The three temporal searches took 145~min in total (42.8, 59.5 and 43.0~min) and the transfer search 21~min. LightGBM's OpenCL GPU backend was 3.0 times slower than the CPU on this workload (47.4~s against 15.8~s for 300 trees on the training split). With 13 features there is too little histogram work per row for the GPU to recover its kernel-launch overhead.
%% GPU timing source: v3_ensemble/README.md ("Note on GPU"); no log file in the repository.

\bibliographystyle{elsarticle-harv}
\bibliography{refs}

\end{document}
